% ==========================================================================
%  SPHERExLens -- presentation specification: every slide, every figure and
%  the complete spoken script. 30 slides; full route 28:10, 15-minute route 18:40.
%  Layout follows the reference "slides and script" document:
%  title page + deck-map table, then one page per slide with
%  "On the slide." / "Figure / diagram." / the figure / "Spoken script."
%  Build:  pdflatex slides_and_script.tex   (twice)
% ==========================================================================
\documentclass[11pt,a4paper]{article}

\usepackage[T1]{fontenc}
\usepackage[utf8]{inputenc}
\usepackage{lmodern}
\usepackage{microtype}
\usepackage[a4paper,margin=2.0cm,top=2.0cm,bottom=2.0cm]{geometry}
\usepackage{amsmath,amssymb}
\usepackage{booktabs,array,longtable,multirow}
\usepackage{graphicx}
\usepackage{xcolor}
\usepackage{enumitem}
\usepackage{titlesec}
\usepackage{fancyhdr}
\usepackage{tikz}
\usetikzlibrary{arrows.meta,positioning,shapes.geometric,calc,fit,backgrounds,patterns,decorations.pathreplacing}
\usepackage{pgfplots}
\pgfplotsset{compat=1.18}
\usepackage{listings}
\usepackage[hidelinks]{hyperref}

% --------------------------------------------------------------- colours ---
\definecolor{ink}{HTML}{14213A}
\definecolor{slate}{HTML}{3F4A5A}
\definecolor{muted}{HTML}{7B8798}
\definecolor{accent}{HTML}{1F6FB2}
\definecolor{accent2}{HTML}{E8F1F8}
\definecolor{good}{HTML}{1E7B4F}
\definecolor{good2}{HTML}{E7F3EC}
\definecolor{bad}{HTML}{B3261E}
\definecolor{bad2}{HTML}{FAECEA}
\definecolor{warn}{HTML}{8A5A00}
\definecolor{warn2}{HTML}{FBF3E2}
\definecolor{rule}{HTML}{C9D0DA}
\definecolor{purple}{HTML}{7C3AED}
\definecolor{purple2}{HTML}{F3E8FF}

% ------------------------------------------------------------ page style ---
\pagestyle{plain}
\setlength{\parindent}{0pt}
\setlength{\parskip}{4pt}

% ---------------------------------------------------------------- macros ---
\newcommand{\arcmin}{\ensuremath{{}^{\prime}}}
\newcommand{\arcsec}{\ensuremath{{}^{\prime\prime}}}
\newcommand{\code}[1]{\texttt{\small\detokenize{#1}}}

% slide header:  Slide 1 --- Title       0:45
% The header is measured first and scaled down to the text width when a long
% title would otherwise run into the running time on the right.
\newsavebox{\slideheadbox}
\newcommand{\slidehead}[3]{%
  \sbox{\slideheadbox}{\large\bfseries Slide #1 --- #2\hspace{1em plus 1fill}#3}%
  \noindent\ifdim\wd\slideheadbox>\linewidth
    \resizebox{\linewidth}{!}{\usebox{\slideheadbox}}%
  \else
    \makebox[\linewidth][s]{\usebox{\slideheadbox}}%
  \fi\par
  \vspace{2pt}\hrule height 0.9pt\vspace{8pt}}

% figure wrapper: sized to fit cleanly without overflowing margins
\newcommand{\figbox}[1]{\begin{center}\resizebox{0.98\linewidth}{!}{#1}\end{center}}


\usepackage[active,tightpage]{preview}
\PreviewEnvironment{tikzpicture}
\begin{document}
\begin{tikzpicture}[font=\footnotesize]
  \foreach \x/\lab in {0/{Visit 1 ($t_1$)}, 4.1/{Visit 2 ($t_2=t_1+\Delta t$)}, 8.2/{Difference ($D$)}, 12.3/{Flaring Monopole}}{
    \draw[rule,thick,fill=white,rounded corners=2pt] (\x,3.0) rectangle (\x+3.8,6.4);
    \node[above=2pt,font=\scriptsize\bfseries,text=ink] at (\x+1.9,6.4) {\lab};
  }
  % Stationary stars
  \foreach \x/\y in {0.5/5.8, 1.2/3.8, 3.2/5.9, 3.3/3.4, 0.9/4.7, 3.2/4.8}{
    \fill[slate!50] (\x,\y) circle (0.06);
    \fill[slate!50] (\x+4.1,\y) circle (0.06);
    \fill[slate!20] (\x+8.2,\y) circle (0.06);
  }
  % Moving source
  \fill[good] (1.4,4.7) circle (0.13);
  \fill[good] (4.1+2.3,4.7) circle (0.13);
  \draw[-{Stealth[length=4pt]},thick,good] (4.1+1.55,4.7) -- (4.1+2.15,4.7);
  \fill[bad] (8.2+1.4,4.7) circle (0.14);
  \fill[good] (8.2+2.3,4.7) circle (0.14);
  \node[anchor=north,font=\tiny,text=muted] at (8.2+1.85,4.45) {dipole ($\mathbf{s}\Delta_{\mathrm{px}}$)};

  % Fading source in Difference panel
  \fill[warn] (2.8,3.9) circle (0.12);
  \fill[bad] (8.2+2.8,3.9) circle (0.13);
  \node[anchor=north,font=\tiny,text=bad] at (8.2+2.8,3.65) {negative lobe};

  % Flaring source in 4th panel
  \fill[warn] (12.3+1.9,4.7) circle (0.24);
  \node[anchor=north,font=\tiny,text=warn] at (12.3+1.9,4.2) {positive monopole};
  \node[anchor=north,font=\tiny,text=muted] at (12.3+1.9,3.75) {flaring accretion / SN};

  % Captions
  \node[anchor=north,font=\tiny,text=muted] at (1.9,2.85) {baseline sky};
  \node[anchor=north,font=\tiny,text=muted] at (4.1+1.9,2.85) {kinematic displacement};
  \node[anchor=north,font=\tiny,text=muted] at (8.2+1.9,2.85) {dipole + fading residual};

  % Systematics warning box
  \node[draw=bad,thick,fill=bad2,rounded corners=2pt,align=left,inner sep=6pt,text width=15.4cm,font=\scriptsize] at (8.05,1.25)
    {\textbf{What survives raw subtraction, and why subtraction alone is not discovery:}\\
     Point-spread function variations between spacecraft roll angles, sub-pixel detector flat-field errors, astrometric jitter, and zodiacal gradients survive subtraction as structured residuals. These instrumental false alarms outnumber genuine astrophysical changes by up to $1000\times$. The scientific task is therefore statistical calibration and multi-wavelength triage.};
\end{tikzpicture}
\end{document}
